\SourceLeaf{328}

Man setze Kürze halber
\[
(p''-p^0)(p'''-p^0) = a
\]
\[
(p'''-p^0)(p'-p^0) = b
\]
\[
(p'-p^0)(p''-p^0) = c
\]
und die Adjuncten dieser Grössen $= A, B, C$.

\Needspace{13\baselineskip}Man setze
\begin{equation}
Q = aB+bC+cA-Ab-Bc-Ca,
\tag*{1)}
\end{equation}
also auch (identisch)
\begin{equation}
Q = (a-b)(A-C)-(A-B)(a-c)
\tag*{2)}
\end{equation}
\begin{equation}
Q = (b-c)(B-A)-(B-C)(b-a)
\tag*{3)}
\end{equation}
\begin{equation}
Q = (c-a)(C-B)-(C-A)(c-b).
\tag*{4)}
\end{equation}

Angewandt auf unsern Fall gibt
\begin{equation}
\begin{aligned}
Q = \;& (p'''-p^0)(P'''-P^0)\{(p''-p^0)(P'-P^0)-(P''-P^0)(p'-p^0)\}\\
+\;& (p'-p^0)(P'-P^0)\{(p'''-p^0)(P''-P^0)-(P'''-P^0)(p''-p^0)\}\\
+\;& (p''-p^0)(P''-P^0)\{(p'-p^0)(P'''-P^0)-(P'-P^0)(p'''-p^0)\},
\end{aligned}
\tag{1}
\end{equation}
\begin{equation}
Q=(p'''-p^0)(p''-p')(P''-P^0)(P'''-P') - (P'''-P^0)(P''-P')(p''-p^0)(p'''-p'),
\tag{2}
\end{equation}
\begin{equation}
Q=(p'-p^0)(p'''-p'')(P'''-P^0)(P'-P'') - (P'-P^0)(P'''-P'')(p'''-p^0)(p'-p''),
\tag{3}
\end{equation}
\begin{equation}
Q=(p''-p^0)(p'-p''')(P'-P^0)(P''-P''') - (P''-P^0)(P'-P''')(p'-p^0)(p''-p'''),
\tag{4}
\end{equation}

woraus alles leicht von selbst folgt.

Offenbar kann man unbeschadet der Allgemeinheit von Anfang an $p^0=0$ setzen, wodurch die Entwickelung noch viel kürzer wird.
